\documentclass[12pt, a4paper]{article}
% ── Encodage & Langue
\usepackage[utf8]{inputenc}
\usepackage[T1]{fontenc}
\usepackage[french]{babel}
% ── Mise en page
\usepackage[top=2.5cm, bottom=2.5cm, left=2.8cm, right=2.8cm]{geometry}
\usepackage{fancyhdr}
\pagestyle{fancy}
\fancyhf{}
\rhead{\small\sffamily G(n,p) — Détection de Fraude \& Cliques}
\lhead{\small\sffamily Algorithmes \& Implémentation Python}
\cfoot{\small\thepage}
\renewcommand{\headrulewidth}{0.4pt}
% ── Typographie
\usepackage{lmodern}
\usepackage{microtype}
\usepackage{setspace}
% ── Couleurs
\usepackage{xcolor}
\definecolor{codebg}    {HTML}{1c2128}
\definecolor{codeframe} {HTML}{30363d}
\definecolor{kw_blue}   {HTML}{79c0ff}
\definecolor{kw_green}  {HTML}{56d364}
\definecolor{kw_yellow} {HTML}{e3b341}
\definecolor{kw_purple} {HTML}{d2a8ff}
\definecolor{kw_cyan}   {HTML}{76e3ea}
\definecolor{kw_red}    {HTML}{ffa198}
\definecolor{txt_white} {HTML}{e6edf3}
\definecolor{txt_muted} {HTML}{8b949e}
\definecolor{algobg}    {HTML}{f0f6ff}
\definecolor{algoborder}{HTML}{0366d6}
\definecolor{accent}    {HTML}{0366d6}
\definecolor{warningbg} {HTML}{fff8e1}
\definecolor{warningborder}{HTML}{f9a825}
% ── Algorithmes
\usepackage[vlined, ruled, linesnumbered]{algorithm2e}
\SetAlgoVlined
\DontPrintSemicolon
% Mots-clés français
\SetKwInput{KwIn}{Entr\'{e}e}
\SetKwInput{KwOut}{Sortie}
\SetKwIF{Si}{SinonSi}{Sinon}{Si}{alors}{Sinon si}{Sinon}{Fin Si}
\SetKwFor{Pour}{Pour}{faire}{Fin Pour}
\SetKwFor{TantQue}{Tant que}{faire}{Fin Tant Que}
\SetKw{Retourner}{retourner}
\SetKw{Et}{et}
\SetKw{Ou}{ou}
\SetKwData{Vrai}{Vrai}
\SetKwData{Faux}{Faux}
\newcommand{\com}[1]{\tcc*[r]{\small\textit{#1}}}
% ── Listings Python
\usepackage{listings}
\lstdefinestyle{python}{
    language         = Python,
    backgroundcolor  = \color{codebg},
    basicstyle       = \footnotesize\ttfamily\color{txt_white},
    keywordstyle     = \color{kw_blue}\bfseries,
    stringstyle      = \color{kw_yellow},
    commentstyle     = \color{kw_green}\itshape,
    numberstyle      = \tiny\color{txt_muted},
    numbers          = left,
    numbersep        = 10pt,
    stepnumber       = 1,
    frame            = single,
    framesep         = 6pt,
    rulecolor        = \color{codeframe},
    breaklines       = true,
    showstringspaces = false,
    tabsize          = 4,
    captionpos       = b,
    xleftmargin      = 18pt,
    xrightmargin     = 4pt,
    morekeywords     = {self, True, False, None, lambda, def, return, if, else, elif, for, in, while, import, from, as, class, try, except},
    emph             = {random, math, sys, list, set, len, max, min, int, str, dict, frozenset, tuple, range, print},
    emphstyle        = \color{kw_cyan},
    literate         =
    {é}{{\'{e}}}1 {è}{{\`{e}}}1 {ê}{{\^{e}}}1
    {à}{{\`{a}}}1 {â}{{\^{a}}}1 {î}{{\^{i}}}1
    {ô}{{\^{o}}}1 {ù}{{\`{u}}}1 {É}{{\'{E}}}1
    {È}{{\`{E}}}1 {ç}{{\c{c}}}1 {—}{{---}}1,
}
\renewcommand{\lstlistingname}{Code Python}
% ── Boîtes tcolorbox
\usepackage{tcolorbox}
\tcbuselibrary{breakable, skins, listings}
\tcbset{
    explication/.style={
        colback     = algobg,
        colframe    = algoborder,
        coltext     = black,
        fonttitle   = \bfseries\sffamily,
        title       = {Explication},
        arc         = 5pt,
        boxrule     = 0.6pt,
        left        = 8pt, right = 8pt,
        top         = 6pt, bottom = 6pt,
        breakable,
    },
    attention/.style={
        colback     = warningbg,
        colframe    = warningborder,
        fonttitle   = \bfseries\sffamily,
        title       = {Remarque},
        arc         = 5pt,
        boxrule     = 0.6pt,
        breakable,
    },
}
% ── Titres
\usepackage{titlesec}
\titleformat{\section}
{\Large\bfseries\sffamily\color{accent}}
{\thesection.}{0.6em}{}
[\vspace{-4pt}\textcolor{accent}{\rule{\textwidth}{0.4pt}}]
\titleformat{\subsection}
{\large\bfseries\sffamily\color{accent!80!black}}
{\thesubsection}{0.5em}{}
\titleformat{\subsubsection}
{\normalsize\bfseries\sffamily\color{black!70}}
{\thesubsubsection}{0.5em}{}
% ── Mathématiques
\usepackage{amsmath, amssymb}
% ── Graphiques
\usepackage{tikz}
\usetikzlibrary{arrows.meta, shapes.geometric, positioning, fit}
% ── Liens
\usepackage{hyperref}
\hypersetup{
    colorlinks  = true,
    linkcolor   = accent,
    urlcolor    = accent,
    pdftitle    = {G(n,p) — Algorithmes de détection de fraude},
    pdfauthor   = {Jeremi Richan Tossy RATODISOA},
}
% ── Table des matières
\usepackage{tocloft}
\renewcommand{\cftsecfont}{\sffamily\bfseries}
\renewcommand{\cftsubsecfont}{\sffamily}

\begin{document}

% ── Page de titre
\begin{titlepage}
    \centering
    \vspace*{1.5cm}
    % Bandeau sombre en haut
    \begin{tikzpicture}
        \fill[color=black!90, rounded corners=8pt]
        (0,0) rectangle (\textwidth, 2.8cm);
        \node[white, font=\Huge\bfseries\sffamily] at (0.5\textwidth, 1.6cm)
        {G(n,\;p) — Détection de Fraude};
        \node[white!70, font=\large\sffamily] at (0.5\textwidth, 0.7cm)
        {Probabilité Adaptative \& Analyse des Cliques};
    \end{tikzpicture}
    
    \vspace{1.5cm}
    
    % Schéma du flux réseau commercial
    \begin{tikzpicture}[
        role/.style={
            circle, draw=algoborder, fill=algobg, line width=1pt,
            minimum size=1.4cm, font=\small\sffamily\bfseries
        },
        susp/.style={
            diamond, draw=warningborder, fill=warningbg, line width=1pt,
            minimum size=1.4cm, font=\small\sffamily\bfseries
        },
        fledge/.style={-Stealth, thick, gray!70, line width=0.8pt},
        fraud/.style={-Stealth, thick, orange!80!red, dashed, line width=1.2pt},
        ]
        \node[role]  (A)  at (0,0)    {\textcolor{green!50!black}{Appro.}};
        \node[role]  (F)  at (3.5,0)  {\textcolor{blue!70!black}{Fourn.}};
        \node[role]  (C)  at (7,0)    {\textcolor{red!70!black}{Cons.}};
        \node[susp]  (S)  at (3.5,-2.5) {\textcolor{orange!70!black}{Susp.}};
        
        \draw[fledge] (A) -- node[above,font=\footnotesize\sffamily,gray]{flux normal} (F);
        \draw[fledge] (F) -- node[above,font=\footnotesize\sffamily,gray]{flux normal} (C);
        \draw[fraud]  (A) to[bend right=25] node[below,font=\footnotesize\sffamily,red]
        {circuit frauduleux} (C);
        \draw[fledge] (S) -- (F);
        \draw[fledge] (S) -- (A);
    \end{tikzpicture}
    
    \vspace{1.5cm}
    
    {\large\sffamily\bfseries Algorithmes couverts}\par
    \vspace{0.4cm}
    \begin{tabular}{cl}
        \textcolor{accent}{\bfseries 1.} & Probabilité adaptative $p(n)$ \\[4pt]
        \textcolor{accent}{\bfseries 2.} & Génération du graphe G(n, p) \\[4pt]
        \textcolor{accent}{\bfseries 3.} & BFS — Parcours en largeur \\[4pt]
        \textcolor{accent}{\bfseries 4.} & Garantie de connexité — Fusion de composantes \\[4pt]
        \textcolor{accent}{\bfseries 5.} & Coloration gloutonne (Welsh-Powell) \\[4pt]
        \textcolor{accent}{\bfseries 6.} & Recherche de Cliques (Bron-Kerbosch avec Pivot) \\[4pt]
        \textcolor{accent}{\bfseries 7.} & Détection de fraude (Rôles \& Structures) \\[4pt]
        \textcolor{accent}{\bfseries 8.} & Disposition aléatoire contrôlée \\
    \end{tabular}
    
    \vfill
    
    {\sffamily\color{gray!60}
        Graphe aléatoire connexe · non orienté · sans boucle \\
        \medskip
        Positions aléatoires · Focus sur les cliques maximales
    }
\end{titlepage}

% ── Table des matières
\tableofcontents
\newpage

\section{Probabilité Adaptative}
\subsection{Explication}
\begin{tcolorbox}[explication]
    La probabilité $p$ n'est pas fixe. Elle est calculée dynamiquement selon la taille $n$ du graphe pour garantir la connexité tout en restant sparse, basée sur le seuil d'Erdős-Rényi $\frac{\ln n}{n}$.
\end{tcolorbox}

\subsection{Algorithme}
\begin{algorithm}[H]
    \caption{Calcul adaptatif de p(n)}
    \KwIn{$n$ : nombre de sommets}
    \KwOut{$p \in (0, 1]$ : probabilité adaptée}
    \Si{$n \leq 1$}{\Retourner $0{,}0$\;}
    \Si{$n \leq 10$}{\Retourner $0{,}60$\;}
    \Si{$n \leq 30$}{\Retourner $0{,}40$\;}
    $\mathit{seuil} \leftarrow \dfrac{\ln n}{n}$\;
    \Si{$n \leq 100$}{\Retourner $\min(1,\; 1{,}5 \times \mathit{seuil})$\;}
    \Si{$n \leq 500$}{\Retourner $\min(1,\; 2{,}0 \times \mathit{seuil})$\;}
    \Retourner $\min(1,\; 2{,}5 \times \mathit{seuil})$\;
\end{algorithm}

\subsection{Code Python}
\begin{lstlisting}[style=python, caption={calculer\_probabilite(n)}]
def calculer_probabilite(n: int) -> float:
    if n <= 1:
        return 0.0
    if n <= 10:
        return 0.60
    if n <= 30:
        return 0.40
    seuil = math.log(n) / n
    if n <= 100:
        return round(min(1.0, 1.5 * seuil), 4)
    if n <= 500:
        return round(min(1.0, 2.0 * seuil), 4)
    return round(min(1.0, 2.5 * seuil), 4)
\end{lstlisting}

\section{Génération du Graphe G(n, p)}
\subsection{Explication}
\begin{tcolorbox}[explication]
    Le modèle \textbf{Erdős–Rényi G(n, p)} génère aléatoirement un graphe non orienté de $n$ nœuds. Pour chaque paire $(i, j)$ avec $i < j$, une arête est ajoutée si un tirage aléatoire est inférieur ou égal à $p$.
\end{tcolorbox}

\subsection{Algorithme}
\begin{algorithm}[H]
    \caption{Génération G(n, p)}
    \KwIn{$n \in \mathbb{N}^*$, $p \in [0,1]$}
    \KwOut{$G$ : graphe non orienté}
    $G \leftarrow$ graphe vide avec les sommets $\{0, \ldots, n-1\}$\;
    \Pour{$i \leftarrow 0$ \textbf{ à } $n - 2$}{
        \Pour{$j \leftarrow i + 1$ \textbf{ à } $n - 1$}{
            \Si{$\mathcal{U}(0,1) \leq p$}{
                Ajouter l'arête $(i, j)$ à $G$\;
            }
        }
    }
    \Retourner $G$\;
\end{algorithm}

\subsection{Code Python}
\begin{lstlisting}[style=python, caption={generer\_graphe(n, p)}]
def generer_graphe(n: int, p: float) -> dict:
    graphe = {i: set() for i in range(n)}
    for i in range(n):
        for j in range(i + 1, n):
            if random.random() <= p:
                graphe[i].add(j)
                graphe[j].add(i)
    return graphe
\end{lstlisting}

\section{BFS — Parcours en Largeur}
\subsection{Explication}
\begin{tcolorbox}[explication]
    Le BFS explore le graphe niveau par niveau depuis un sommet de départ. Il est utilisé ici pour identifier les composantes connexes.
\end{tcolorbox}

\subsection{Algorithme}
\begin{algorithm}[H]
    \caption{BFS — Parcours en largeur}
    \KwIn{$G=(V,E)$, $s$ : sommet de départ}
    \KwOut{$\mathit{visites}$ : ensemble des sommets atteignables}
    $\mathit{visites} \leftarrow \{s\}$ ; $\mathit{file} \leftarrow [s]$ ; $\mathit{tete} \leftarrow 0$\;
    \TantQue{$\mathit{tete} < |\mathit{file}|$}{
        $v \leftarrow \mathit{file}[\mathit{tete}]$ ; $\mathit{tete} \leftarrow \mathit{tete} + 1$\;
        \Pour{chaque voisin $w$ de $v$}{
            \Si{$w \notin \mathit{visites}$}{
                $\mathit{visites} \leftarrow \mathit{visites} \cup \{w\}$ ;
                $\mathit{file} \leftarrow \mathit{file} + [w]$\;
            }
        }
    }
    \Retourner $\mathit{visites}$\;
\end{algorithm}

\subsection{Code Python}
\begin{lstlisting}[style=python, caption={bfs(graphe, depart)}]
def bfs(graphe: dict, depart: int) -> set:
    visites = {depart}
    file = [depart]
    tete = 0
    while tete < len(file):
        sommet = file[tete]
        tete += 1
        for voisin in graphe[sommet]:
            if voisin not in visites:
                visites.add(voisin)
                file.append(voisin)
    return visites
\end{lstlisting}

\section{Connexité — Garantie et Composantes}
\subsection{Explication}
\begin{tcolorbox}[explication]
    L'algorithme identifie toutes les composantes connexes via des BFS successifs, puis relie les composantes deux à deux par une arête aléatoire jusqu'à obtenir un graphe unique connexe.
\end{tcolorbox}

\subsection{Algorithme}
\begin{algorithm}[H]
    \caption{Garantie de connexité}
    \KwIn{$G$ : graphe potentiellement non connexe}
    \KwOut{$G$ : graphe connexe}
    $\mathit{composantes} \leftarrow$ TrouverComposantesConnexes($G$)\;
    \TantQue{$|\mathit{composantes}| > 1$}{
        $c_1 \leftarrow \mathit{composantes}[0]$ ; $c_2 \leftarrow \mathit{composantes}[1]$\;
        $u \leftarrow$ \textsc{ChoixAléatoire}$(c_1)$\;
        $v \leftarrow$ \textsc{ChoixAléatoire}$(c_2)$\;
        Ajouter l'arête $(u, v)$ à $G$\;
        $\mathit{composantes} \leftarrow$ TrouverComposantesConnexes($G$)\;
    }
    \Retourner $G$\;
\end{algorithm}

\subsection{Code Python}
\begin{lstlisting}[style=python, caption={rendre\_connexe(graphe)}]
def rendre_connexe(graphe: dict) -> None:
    composantes = composantes_connexes(graphe)
    while len(composantes) > 1:
        s1 = random.choice(list(composantes[0]))
        s2 = random.choice(list(composantes[1]))
        graphe[s1].add(s2)
        graphe[s2].add(s1)
        composantes = composantes_connexes(graphe)
\end{lstlisting}

\section{Coloration Gloutonne — Welsh-Powell}
\subsection{Explication}
\begin{tcolorbox}[explication]
    La coloration gloutonne assigne à chaque sommet la plus petite couleur non utilisée par ses voisins. Les sommets sont triés par degré décroissant (Welsh-Powell) pour optimiser le nombre de couleurs.
\end{tcolorbox}

\subsection{Algorithme}
\begin{algorithm}[H]
    \caption{Coloration gloutonne Welsh-Powell}
    \KwIn{$G=(V,E)$}
    \KwOut{$\mathit{couleur}[v]$ pour tout $v \in V$}
    $\mathit{ordre} \leftarrow$ trier $V$ par degré décroissant\;
    $\mathit{couleur} \leftarrow \{\}$\;
    \Pour{chaque sommet $v$ dans $\mathit{ordre}$}{
        $C_{\text{voisins}} \leftarrow \{ \mathit{couleur}[w] \mid w \in N(v), w \in \mathit{couleur} \}$\;
        $c \leftarrow 0$\;
        \TantQue{$c \in C_{\text{voisins}}$}{$c \leftarrow c + 1$\;}
        $\mathit{couleur}[v] \leftarrow c$\;
    }
    \Retourner $\mathit{couleur}$\;
\end{algorithm}

\subsection{Code Python}
\begin{lstlisting}[style=python, caption={colorier(graphe)}]
def colorier(graphe: dict) -> dict:
    ordre = sorted(graphe, key=lambda s: len(graphe[s]), reverse=True)
    couleur = {}
    for sommet in ordre:
        couleurs_voisins = {couleur[v] for v in graphe[sommet] if v in couleur}
        c = 0
        while c in couleurs_voisins:
            c += 1
        couleur[sommet] = c
    return couleur
\end{lstlisting}

\section{Cliques — Algorithme de Bron-Kerbosch}
\subsection{Explication}
\begin{tcolorbox}[explication]
    L'algorithme de Bron-Kerbosch énumère toutes les cliques maximales par backtracking. Une optimisation par \textbf{pivot} est utilisée pour réduire le nombre d'appels récursifs. Seules les cliques de taille $\geq 3$ sont retenues.
\end{tcolorbox}

\subsection{Algorithme}
\begin{algorithm}[H]
    \caption{Bron-Kerbosch avec Pivot}
    \KwIn{$G$, $R$ (résultat), $P$ (candidats), $X$ (exclus)}
    \KwOut{Liste des cliques maximales}
    \Si{$P = \varnothing$ \Et $X = \varnothing$}{
        \Si{$|R| \geq 3$}{Ajouter $R$ aux résultats}
        \Retourner
    }
    $u \leftarrow$ pivot dans $P \cup X$ maximisant $|P \cap N(u)|$\;
    \Pour{chaque sommet $v$ dans $P \setminus N(u)$}{
        BronKerbosch($G$, $R \cup \{v\}$, $P \cap N(v)$, $X \cap N(v)$)\;
        $P \leftarrow P \setminus \{v\}$\;
        $X \leftarrow X \cup \{v\}$\;
    }
\end{algorithm}

\subsection{Code Python}
\begin{lstlisting}[style=python, caption={bron\_kerbosch(...) et trouver\_cliques(...)}]
def bron_kerbosch(graphe, R, P, X, cliques):
    if not P and not X:
        if len(R) >= 3:
            cliques.append(frozenset(R))
        return
    
    # Choix du pivot
    pivot = max(P | X, key=lambda u: len(graphe[u] & P))
    
    for sommet in list(P - graphe[pivot]):
        voisins = graphe[sommet]
        bron_kerbosch(graphe, R | {sommet}, 
                      P & voisins, X & voisins, cliques)
        P = P - {sommet}
        X = X | {sommet}

def trouver_cliques(graphe: dict) -> list:
    cliques = []
    sommets = set(graphe.keys())
    bron_kerbosch(graphe, set(), sommets, set(), cliques)
    return cliques
\end{lstlisting}

\section{Rôles et Détection de Fraude}
\subsection{Explication}
\begin{tcolorbox}[explication]
    Les rôles sont assignés selon la couleur (plafonnée à 3). La fraude est détectée si :
    \begin{enumerate}
        \item Une arête relie directement un Consommateur (Rôle 0) et un Approvisionneur (Rôle 2).
        \item Un nœud a un rôle Suspect (Rôle 3, issu d'une coloration élevée).
    \end{enumerate}
\end{tcolorbox}

\subsection{Algorithme}
\begin{algorithm}[H]
    \caption{Détection des circuits frauduleux}
    \KwIn{$G=(V,E)$, $\mathit{roles}[v]$}
    \KwOut{$S_{\text{fraude}} \subseteq V$, $A_{\text{fraude}} \subseteq E$}
    $S_{\text{fraude}} \leftarrow \varnothing$ ; $A_{\text{fraude}} \leftarrow \varnothing$\;
    \Pour{chaque arête $(u,v) \in E$ avec $u < v$}{
        \Si{$(\mathit{roles}[u]=0 \;\Et\; \mathit{roles}[v]=2)$ \Ou $(\mathit{roles}[u]=2 \;\Et\; \mathit{roles}[v]=0)$}{
            $A_{\text{fraude}} \leftarrow A_{\text{fraude}} \cup \{(u,v)\}$\;
            $S_{\text{fraude}} \leftarrow S_{\text{fraude}} \cup \{u, v\}$\;
        }
    }
    \Pour{chaque sommet $v \in V$}{
        \Si{$\mathit{roles}[v] = 3$}{
            $S_{\text{fraude}} \leftarrow S_{\text{fraude}} \cup \{v\}$\;
        }
    }
    \Retourner $S_{\text{fraude}},\; A_{\text{fraude}}$\;
\end{algorithm}

\subsection{Code Python}
\begin{lstlisting}[style=python, caption={detecter\_fraude(graphe, roles)}]
def detecter_fraude(graphe: dict, roles: dict) -> tuple:
    sommets_frauduleux = set()
    aretes_frauduleuses = set()
    
    for u in graphe:
        for v in graphe[u]:
            if u < v:
                ru, rv = roles[u], roles[v]
                if (ru == 0 and rv == 2) or (ru == 2 and rv == 0):
                    aretes_frauduleuses.add((u, v))
                    sommets_frauduleux.update([u, v])

    for s, role in roles.items():
        if role == 3:
            sommets_frauduleux.add(s)

    return sommets_frauduleux, aretes_frauduleuses
\end{lstlisting}

\section{Disposition Aléatoire}
\subsection{Explication}
\begin{tcolorbox}[explication]
    Les nœuds sont placés aléatoirement dans un rectangle défini. Une tentative de contrôle de distance minimale est effectuée pour éviter la superposition totale des nœuds, bien que cela reste une disposition force-directed simplifiée.
\end{tcolorbox}

\subsection{Code Python}
\begin{lstlisting}[style=python, caption={disposition\_aleatoire(graphe)}]
def disposition_aleatoire(graphe, largeur=2.0, hauteur=2.0, dist_min=0.08):
    positions = {}
    demi_l = largeur / 2
    demi_h = hauteur / 2
    
    for sommet in graphe:
        meilleure_pos = None
        meilleure_dist = -1.0
        
        for _ in range(30): # Max tentatives
            x = random.uniform(-demi_l, demi_l)
            y = random.uniform(-demi_h, demi_h)
            
            if not positions:
                meilleure_pos = (x, y)
                break
            
            # Calcul distance min aux autres
            dist_min_loc = min(
                ((x - px)**2 + (y - py)**2)**0.5 
                for px, py in positions.values()
            )
            
            if dist_min_loc >= dist_min:
                meilleure_pos = (x, y)
                break
                
            if dist_min_loc > meilleure_dist:
                meilleure_dist = dist_min_loc
                meilleure_pos = (x, y)
                
        positions[sommet] = meilleure_pos
    return positions
\end{lstlisting}
\end{document}
